\documentclass[conference]{IEEEtran}
\IEEEoverridecommandlockouts
\usepackage{cite}
\usepackage{amsmath,amssymb,amsfonts}
\usepackage{graphicx}
\usepackage{booktabs}
\usepackage{url}
\usepackage[hidelinks]{hyperref}
\usepackage{xcolor}

\begin{document}

\title{ExpenseLM: Distilling Policy-Aware Expense Extraction\\into a 4B Language Model}

\author{\IEEEauthorblockN{Suyash Kumar Singh (Roll No. 23035010145)}
\IEEEauthorblockA{Term Project Report, Trimester 8\\
B.Sc. (Honours) Data Science and Artificial Intelligence\\
Indian Institute of Technology Guwahati, India\\
s.suyash@op.iitg.ac.in}}

\maketitle

\begin{abstract}
Converting messy, real-world expense text---chat messages, receipt OCR dumps,
forwarded emails---into clean, policy-checked structured records is a common
business task, today handled either manually or by calling expensive
frontier-model APIs on every message. We test whether a small ($\sim$4-billion
parameter) open model, fine-tuned on a modest task-specific dataset, can
approach frontier accuracy on this task at near-zero inference cost. We build
a complete post-training pipeline: a procedurally generated, policy-labelled
dataset; supervised fine-tuning (SFT) of Qwen3-4B with QLoRA; and a custom
evaluation harness measuring schema validity, field accuracy, expense-level
$F_1$, and policy-reasoning accuracy. Fine-tuning more than doubled
policy-reasoning accuracy over the base model (24.7\% to 55.5\%) and lifted
relative-date resolution from 43\% to 96\%, closing roughly half the gap to a
frontier ceiling using 1{,}600 training examples and no paid training compute.
We additionally isolate and fix a concrete decoding failure mode. The study
demonstrates, end to end, how a narrow capability can be distilled from a large
model into a cheap, private, self-hosted one, and quantified rigorously.
\end{abstract}

\begin{IEEEkeywords}
large language models, fine-tuning, LoRA, QLoRA, structured extraction,
model distillation, evaluation.
\end{IEEEkeywords}

\section{Introduction}
Businesses receive expense information as unstructured text: a Hinglish chat
message (``bhai 450 ka cab liya airport tak''), a broken receipt OCR dump, or a
forwarded email with signatures and quoted threads as noise. Turning this into a
clean record---amount, currency, category, merchant, date, reimbursability, and
policy flags---is done either by hand or by sending every message to a
frontier-model API, which is costly and moves private financial data off-site.

This mirrors one of the most common production patterns in modern applied AI:
\emph{distilling a narrow capability from a large ``teacher'' model into a
small, cheap, self-hosted ``student'' model} \cite{deepseek}. The task is
non-trivial because the model must not merely extract fields but \emph{apply a
policy}---per-category spending limits, receipt thresholds, and reimbursability
rules---which requires reading comprehension and reasoning rather than pattern
matching.

\subsection*{Objectives}
\begin{itemize}
  \item Specify the extraction task with a strict output schema and an
        explicit expense policy that the model must apply.
  \item Construct a diverse, correctly labelled dataset covering chat, OCR,
        email, and multilingual (English/Hinglish) inputs.
  \item Fine-tune a 4B open model using parameter-efficient QLoRA and compare
        it against zero-shot, few-shot, and a frontier-model ceiling.
  \item Build a rigorous, self-written evaluation harness and quantify what
        each stage of training buys.
  \item Analyse residual failures to characterise \emph{when} the model can be
        trusted.
\end{itemize}

\section{Related Work}
Parameter-efficient fine-tuning makes adapting large models tractable on
commodity hardware. \textbf{LoRA} \cite{lora} freezes the base weights and learns
a low-rank update; \textbf{QLoRA} \cite{qlora} adds 4-bit quantisation of the
frozen base, enabling billion-parameter fine-tuning in a few gigabytes of GPU
memory. \textbf{Direct Preference Optimization (DPO)} \cite{dpo} aligns a model
to preferred outputs without a separate reward model. At industrial scale,
\textbf{DeepSeek-R1} \cite{deepseek} distils reasoning from a large model into
small dense models by supervised fine-tuning on teacher-generated traces---the
same teacher-to-student pattern this project applies at a single-GPU scale. We
use \textbf{Qwen3-4B-Instruct} \cite{qwen} as the base model and the
\textbf{Unsloth} \cite{unsloth} and \textbf{TRL} \cite{trl} libraries for
efficient training.

\section{Methodology}

\subsection{Task Specification}
The input is a raw text string; the output is a JSON object with a list of
expense objects and an overall confidence label. Each expense has the fields
\texttt{amount}, \texttt{currency} (ISO 4217 or \texttt{UNKNOWN}),
\texttt{category} (a closed enum), \texttt{merchant}, \texttt{date}
(resolved against a reference date supplied in the prompt),
\texttt{description}, \texttt{reimbursable}, and \texttt{policy\_flags} (a
closed set such as \texttt{over\_limit}, \texttt{missing\_receipt}). At
inference the system prompt carries the schema, a reference date, and a compact
policy block; this same prompt is used for \emph{every} system so that measured
differences are attributable to the training method, not prompt wording.

\subsection{Dataset Construction}
Because manually labelling thousands of examples is infeasible, we generate
data \emph{procedurally}: each example is composed from orthogonal axes
(source type $\times$ language $\times$ ambiguity $\times$ number of expenses
$\times$ policy situation), and the gold label is \emph{computed from the
policy rules in code}. Consequently, input and label can never disagree---the
label is a computation, not an opinion. A programmatic perturbation pass then
injects realism (typos, OCR glyph confusions such as $0\!\leftrightarrow\!O$,
and currency-symbol variants) to the \emph{input only}, preserving the label.

A contamination check using token-set fuzzy matching removes any training
example that is a near-duplicate of a test example; on the first pass it culled
1{,}134 template twins, precisely the leakage it exists to prevent. The final
splits are \textbf{1{,}600 train / 100 dev / 200 test}, all schema-valid, with
zero train--test contamination. During validation, a disagreement with a
frontier model surfaced 349 labels naming a merchant absent from the input;
these were repaired deterministically. The test set was frozen and used exactly
once per system.

\subsection{Model and Fine-Tuning}
We fine-tune with QLoRA. The frozen base weights $W$ are stored in 4-bit NF4;
a low-rank adapter is learned such that the effective weight is
\begin{equation}
W' = W + \frac{\alpha}{r} B A, \quad A \in \mathbb{R}^{r \times d},\;
B \in \mathbb{R}^{d \times r},
\end{equation}
with rank $r=16$, scaling $\alpha=16$, applied to all seven linear projections
of the transformer. Only $A$ and $B$ are trained (under 1\% of parameters).
Training used learning rate $2\times10^{-4}$, a cosine schedule, two epochs,
effective batch size 16 via gradient accumulation, 8-bit AdamW, and gradient
checkpointing, on a single free-tier T4 GPU. Validation loss fell from
$0.112$ to $0.063$ with no over-fitting.

\subsection{Evaluation Harness}
Evaluation is deliberately hand-written ($\sim$300 lines) rather than delegated
to a framework, and is built \emph{before} any training so that metric
definitions cannot be tuned to flatter the model. Five metric layers each
isolate a distinct question:
\begin{itemize}
  \item \textbf{Schema validity}: does the output parse and satisfy the strict
        \texttt{pydantic} contract?
  \item \textbf{Field accuracy}: per-field correctness (amounts within
        $\pm 0.01$).
  \item \textbf{Expense-level $F_1$}: for multi-expense inputs, are all
        expenses found with no hallucinated extras?
  \item \textbf{Policy accuracy}: are \texttt{reimbursable} and the exact set
        of \texttt{policy\_flags} both correct---the reasoning metric.
  \item \textbf{Cost/latency}: local throughput versus API price.
\end{itemize}
A key sub-problem is \emph{expense alignment}: predicted expenses may be
reordered, so before grading fields we match predictions to gold greedily on
(amount within tolerance, then category). The harness carries 29 unit tests and
a gold-echo integration test (a system returning gold scores exactly $1.0$),
validating the grader itself.

\section{Experimental Setup}
All systems are evaluated on the identical 200-example test set:
\begin{itemize}
  \item \textbf{E0}: Qwen3-4B-Instruct, zero-shot.
  \item \textbf{E1}: Qwen3-4B-Instruct, 5-shot.
  \item \textbf{E2}: E0 $+$ SFT (QLoRA).
  \item \textbf{E4}: a frontier model, zero-shot---the accuracy ceiling.
\end{itemize}
Preference tuning (E3, DPO) and quantised GGUF export (E5) are implemented but
were not executed to completion due to exhaustion of free-tier GPU quotas; they
are discussed as future scope.

\section{Results and Discussion}

\begin{table}[t]
\caption{Results on the 200-example test set. E2$^{\dagger}$ is depressed by a
decoding artifact (Sec.~\ref{sec:artifact}); the true capability is in
Table~\ref{tab:cond}.}
\label{tab:main}
\centering
\small
\begin{tabular}{lcccc}
\toprule
System & Schema & Exp. $F_1$ & Date & Policy \\
\midrule
E0 zero-shot        & 99.5\% & 0.980 & 42.9\% & 27.1\% \\
E1 5-shot           & 98.0\% & 0.988 & 79.5\% & 34.3\% \\
E2 SFT$^{\dagger}$  & 81.0\% & 0.898 & 77.9\% & 45.2\% \\
E4 frontier         & 100\%  & 1.000 & 100\%  & 90.4\% \\
\bottomrule
\end{tabular}
\end{table}

\subsection{A Decoding Artifact, Isolated}
\label{sec:artifact}
On evaluation, E2 produced valid JSON on 162 of 200 examples; the other 38 are
an identical failure---the model emits a \texttt{<tool\_call>} token
repeatedly and never answers. The cause is that the base model's chat template
exposes a tool-calling token that greedy decoding occasionally collapses into.
This is a \emph{decoding} artifact, not a learning failure, and is removed by
suppressing that token at generation. Because every failed example also scores
zero on every field, this single artifact depresses all of E2's aggregate
numbers in Table~\ref{tab:main}.

\subsection{The True Effect of Fine-Tuning}
Restricting all systems to the 162 examples E2 answered isolates learned
capability from the decoding bug (Table~\ref{tab:cond}).

\begin{table}[t]
\caption{Apples-to-apples comparison on the same 162 examples.}
\label{tab:cond}
\centering
\small
\begin{tabular}{lcccc}
\toprule
System & Exp. $F_1$ & Policy & Date & Category \\
\midrule
E0 zero-shot & 0.980 & 24.7\% & 43\% & --- \\
E1 5-shot    & 0.986 & 32.4\% & 77\% & --- \\
\textbf{E2 SFT} & \textbf{1.000} & \textbf{55.5\%} & \textbf{96\%} & \textbf{99.6\%} \\
E4 frontier  & 1.000 & 89.9\% & 100\% & 100\% \\
\bottomrule
\end{tabular}
\end{table}

Four observations follow. \textbf{(1) Extraction is easy; reasoning is hard.}
Even zero-shot, the base model emits valid JSON 99.5\% of the time and finds the
right expenses ($F_1=0.98$); what it cannot do is apply the policy (27\%) or
resolve relative dates (43\%). \textbf{(2) Few-shot teaches format, not
reasoning.} Five in-context examples lifted date resolution ($43\%\!\to\!79\%$)
but policy only marginally ($27\%\!\to\!34\%$). \textbf{(3) SFT teaches the
reasoning.} On the shared set, fine-tuning more than doubled policy accuracy
($24.7\%\!\to\!55.5\%$), nearly solved date resolution ($43\%\!\to\!96\%$), and
made expense-finding perfect---the central result of this work. \textbf{(4) The
model sits between few-shot and the frontier}, closing roughly half the policy
gap to a frontier model using a 4B model, 1{,}600 synthetic examples, and no
paid training compute. Once deployed on local hardware the marginal inference
cost is effectively zero, versus a per-message API charge for the frontier
system.

\subsection{Failure Analysis}
Aside from the decoding artifact, residual errors among answered examples
concentrate in policy edge cases---exactly the 55.5\%-versus-90\% gap that
preference tuning (DPO) is designed to close by training on the model's own
harvested mistakes. This tells a practitioner \emph{when} to trust the model
(clean single-expense inputs) and when to route to a human (ambiguous
reimbursability), which is the practical value of rigorous evaluation.

\section{Conclusion and Future Scope}
We showed, end to end, that supervised fine-tuning of a 4B open model teaches
the policy reasoning and date arithmetic that prompting alone cannot, more than
doubling policy accuracy over the base and closing about half the distance to a
frontier ceiling---at near-zero inference cost and full data privacy. Equally
important, a purpose-built evaluation harness made the result legible and
exposed a concrete, fixable decoding failure that aggregate accuracy alone
would have hidden.

\textbf{Future scope.} (i) \emph{Preference tuning (DPO)}: harvest the SFT
model's real mistakes into (chosen, rejected) pairs and train with DPO to close
the remaining policy gap; the pipeline is implemented. (ii) \emph{Quantised
deployment}: export a 4-bit GGUF and measure the accuracy cost of quantisation
and the local-versus-API economics. (iii) \emph{Data diversity}: blend
frontier-generated examples into the procedural corpus to reduce template bias.
(iv) \emph{Reinforcement learning with a verifiable reward}: because our grader
is a programmatic reward, GRPO-style training against it is a natural extension.

\begin{thebibliography}{9}
\bibitem{lora} E.~J.~Hu \emph{et al.}, ``LoRA: Low-Rank Adaptation of Large
Language Models,'' \emph{Int. Conf. Learning Representations (ICLR)}, 2022.
\bibitem{qlora} T.~Dettmers, A.~Pagnoni, A.~Holtzman, and L.~Zettlemoyer,
``QLoRA: Efficient Finetuning of Quantized LLMs,'' \emph{Advances in Neural
Information Processing Systems (NeurIPS)}, 2023.
\bibitem{dpo} R.~Rafailov \emph{et al.}, ``Direct Preference Optimization: Your
Language Model is Secretly a Reward Model,'' \emph{NeurIPS}, 2023.
\bibitem{deepseek} DeepSeek-AI, ``DeepSeek-R1: Incentivizing Reasoning
Capability in LLMs via Reinforcement Learning,'' arXiv:2501.12948, 2025.
\bibitem{qwen} Qwen Team, ``Qwen3 Technical Report,'' arXiv, 2025.
\bibitem{unsloth} Unsloth AI, ``Unsloth: Fast and memory-efficient LLM
fine-tuning,'' \url{https://github.com/unslothai/unsloth}, 2024.
\bibitem{trl} L.~von~Werra \emph{et al.}, ``TRL: Transformer Reinforcement
Learning,'' \url{https://github.com/huggingface/trl}, 2020.
\end{thebibliography}

\end{document}
